% !TeX root = main.tex

% 단일 대회에서는 이 파일에 출제진 후기를 작성한다.
% \authorreview{아이디}{후기 본문}을 필요한 만큼 추가한다.
\newcommand{\authorreview}[2]{%
    \par\noindent
    {\color{solutio-purple}\bfseries #1}: #2\par
    \medskip
}

\section{출제진 후기}

% 후기가 길어지면 frame을 추가하고 작성란을 나눈다.
% label=afterword는 첫 후기 프레임에만 둔다.
\begin{frame}[label=afterword]{출제진 후기}
     \authorreview{yhjang04}{
        \textcolor{solutio-gray}{
            이번 KGUPC 다들 많은 관심 가져주셔서 정말 감사드립니다.
        }
    }

    \authorreview{areian13}{
        \textcolor{solutio-gray}{
            후기를 작성해 주세요.
        }
    }

    \authorreview{ehdnjs14}{
        \textcolor{solutio-gray}{
            후기를 작성해 주세요.
        }
    }

    \authorreview{mingmingmon}{
        \textcolor{solutio-gray}{
            제 일상을 녹인 문제들을 출제해서 스토리 짜는데 재밌었습니다. \\
            1회랑 다르게 타대학에서 참가자분들이 많은 관심을 주셔서 대회가 더 풍성해 진 것 같아요. 다음에는 귀여운 특별상도 사비로 기획해오겠습니다 (제 문제 위주로 ㅎㅎ) 
        }
    }


\end{frame}
